\subsection{A marked square and the determinant increment [D]}
It remains to estimate the increment of $R_{\rm det}$. We can measure the effect of $e=(r,c)$
on $L_T$ by varying its weight or by comparing the board with the minor left after occupying it.
For an allowed $e$, let $\pi^{\rm der}_e=\partial_{\log w_e}L_T=p_e(1+\hat c_e)$,
$\varepsilon_e=p_e/P_e-1$, and $r_e=L_T(M_e)-L_T(M)-\log(\pi^{\rm der}_e/w_e)$.
Theorem A for rooks gives $|\hat c_e|\le C/k$ and $|r_e|\le C/k^2$. Thus $r_e$ measures the
discrepancy between the two operations through second order. Using the scaling of the minor
$P'=\operatorname{diag}(x')P^{(e)}\operatorname{diag}(y')$, where $P^{(e)}$ is the submatrix
of $P$, gives the exact identity
\begin{align*}
 R_{\rm det}(M_e)-R_{\rm det}(M)&=-r_e-\mathsf E_e,\\
 \mathsf E_e={}&\log(1+\varepsilon_e)+\log(1+\hat c_e)
          +\sum_{i\ne r}\log x'_i+\sum_{l\ne c}\log y'_l\\
 &+(k-1)\log(1-1/k)
  +\tfrac12\log\frac{\det(I-b_{P'}^\top b_{P'})}{\det(I-b_P^\top b_P)}\\
 &+\frac{\delta_1((k-1)P')}{k-1}-\frac{\delta_1(kP)}k .
\end{align*}
Indeed, $\Delta=-L_T(P)+\log(k!/k^k)-\frac12\log\det(I-b_P^\top b_P)$.
The covariance of $L_T$ under scaling gives the two sums of potentials. The sign of the size term
comes from the ratio of the two factorials. To prove $(A^+)$ we still need to gain one power of
$k$ in the signed sum $r_e+\mathsf E_e$. The exact expression for $L_T$ retains the quadratic
displacement of the stationary point, usually denoted by $\zeta$. Bounding the eleven second-order
remainders in absolute value would lose their cancellation.

\paragraph{Independent parameters.}
To retain the signed sum, we place both boards in the same family. We vary separately the mass of
row $r$, the mass of column $c$, the inverse size, and the weight of square $e$. Fix $W=kP$ and
take each row and column as a type. With $u=1/k$, let $A_-,B_-$ be the uniform probability
measures that omit the distinguished row and column, respectively. Set
\[
  A_s=(1-s)A_-+s\delta_r,\qquad B_t=(1-t)B_-+t\delta_c,\qquad h=1/J .
\]
Multiply only $W_{rc}$ by $v$. With respect to these measures, the scaled and stationary kernels
satisfy $N=W^v\exp(a_i+b_l)$ and $q(1-hq)=W^v\exp(z_i+y_l)$ with unit marginals.
With the gauge fixed, the inverse of the derivative is uniformly bounded in supremum norm.
The gap controls the quadratic norm, and the integral part maps the quadratic norm to the
supremum norm because $W\le M_0$. Perturbing the measures costs $O(|s-u|+|t-u|)$, and the
marked weight costs $O(\max(|s|,|t|))$ for $|v|<2$. Iteration with the inverse at the center
converges on fixed disks for $h,s,t$ depending only on the class and gives the holomorphic branch.
The determinants of the centered kernels $V$ that occur take the form
\[
 D(V)=\sum_{\ell\ge0}\frac{(-1)^\ell}{(\ell!)^2}
  \int\det[V(i_a,l_b)]_{a,b\le\ell}^{\,2}\,
               dA_s^{\otimes\ell}dB_t^{\otimes\ell}.
\]
Cauchy--Binet and Hadamard give the bound $\ell^\ell C^\ell/(\ell!)^2$. The gap and the
Hilbert--Schmidt norm at the center give a positive lower bound for the determinant there.
The series and Cauchy's estimates let us take its logarithm on uniformly smaller disks. In the
expression for $\Delta$ we replace $\log(J!/J^J)-\frac12\log(4\pi J)$ with
$-1/h-\frac12\log2+h/12$ and keep the Stirling remainder $\eta_J=O(J^{-3})$ separate.
For the resulting expression $D_*$, the product $h^2D_*$ is bounded and holomorphic on these disks.
For positive rational masses, replicating each type $J A_s(i)$ or $J B_t(l)$ times preserves
the operator on types and adds only zero singular values. Applying the deterministic part of G1
to these replicas gives $D_*=\delta_1/J+O(J^{-2})$. Thus the first three coefficients of $h^2D_*$
vanish and the coefficient of degree three is $\delta_1$. By continuity and the holomorphic
identity theorem, Cauchy's estimates give $R_{\rm det}=\eta_J+h^2D_2(h,s,t,v)$, with $D_2$ and
its first derivatives bounded on smaller disks. On a face of zero mass, only the visible part is
replicated. Its gauge drops out of the sums of potentials, and its own $\delta_1$ is identified.

Passing to the minor replaces $(h,s,t)=(u,u,u)$ by $(u/(1-u),0,0)$. The change in $h$ is
$O(u^2)$ and the change in each mass is $O(u)$, giving $\Delta_e R_{\rm det}=O(u^3)$.
The identity above proves $(A^+)$ with the stated sign. On the face $s=0$ or $t=0$, the
integrated expressions are independent of $v$. Thus $\partial_vD_2=stH$ with $H$ bounded by
Cauchy's estimates. At $s=t=u$ the derivative of $R_{\rm det}$ with respect to $v$ is $O(u^4)$
all the way to $v=0$. Integration over $[0,1]$ gives the deletion estimate in $(D^+)$. Adding
the increments from $(L^{*+})$ gives $(T^+)$ for the total remainder $R$. Thus we can compare
the approximation for the same pair before and after passing to a further minor that retains it.
